\section{Extracting state probabilities from a generating function}
\label{m:app:extract}

Every generating function obtained above arrived as a formula rather than as an
explicit series --- the absorption-only case, for instance, as
\begin{equation}
\label{m:eq:examplePgf}
G(x,y,t) = \Bigl(\bigl(1-\ee^{-\alpha t}\bigr)x + \ee^{-\alpha t}y\Bigr)^n .
\end{equation}
Such a formula can always be converted back to the series form
\begin{equation}
\label{m:eq:series1}
G(x,y,t) = \sum_{i=0}^{\infty}\sum_{j=0}^{\infty}p_{i,j}(t)\,x^iy^j ,
\end{equation}
which recovers the time-dependent state probabilities. All that is required is that
$G$ be repeatedly partially differentiable in $x$ and $y$.

The reason is the bivariate Maclaurin expansion,
\begin{equation*}
f(x,y) = \sum_{k=0}^{\infty}\frac{1}{k!}\sum_{\varrho=0}^{k}\binom{k}{\varrho}
\partial_x^{\varrho}\partial_y^{\,k-\varrho}f\Big|_{(0,0)}\;x^{\varrho}y^{k-\varrho},
\end{equation*}
whose coefficient of $x^iy^j$, taking $k=i+j$ and $\varrho=i$, is
$\frac{1}{(i+j)!}\binom{i+j}{i}\partial_x^i\partial_y^{\,j}f\big|_{(0,0)}$. Hence
\begin{equation}
\label{m:eq:stateProb}
p_{i,j}(t) = \frac{1}{i!\,j!}\,\partial_x^i\partial_y^{\,j}G\Big|_{(0,0)} .
\end{equation}
Equivalently, and more stably in floating point when $G$ is a special-function
expression whose high derivatives are unwieldy, the same coefficient is the
bivariate Cauchy integral
\begin{equation}
\label{m:eq:bivariateCauchy}
p_{i,j}(t) = \frac{1}{(2\pi\mathrm i)^2}
\oint_{|x|=\varrho_x}\oint_{|y|=\varrho_y}
\frac{G(x,y,t)}{x^{\,i+1}y^{\,j+1}}\,\dd y\,\dd x
\end{equation}
over circles inside the domain of analyticity.

\subsection*{The multilinear case}

Both \cref{m:eq:Gabsonly} and \cref{m:eq:Gabsdeath} have the form
\begin{equation}
\label{m:eq:multilinear}
G(x,y,t) = \bigl(c(t) + f(t)\,x + g(t)\,y\bigr)^n ,
\end{equation}
for which the derivatives are immediate. Each differentiation lowers the power by
one and brings down a factor, so
\begin{equation}
\label{m:eq:derivs}
\partial_x^{a}\partial_y^{\,b}G
= \frac{n!}{(n-a-b)!}\,f^{a}g^{b}\bigl(c+fx+gy\bigr)^{n-a-b}
\qquad(a+b\le n),
\end{equation}
and zero for $a+b>n$. Evaluating at the origin and applying \cref{m:eq:stateProb},
\begin{equation}
\label{m:eq:stateProbSpec}
p_{i,j}(t) = \frac{n!}{i!\,j!\,(n-i-j)!}\;f(t)^{i}g(t)^{j}c(t)^{\,n-i-j},
\qquad i+j\le n,
\end{equation}
and $p_{i,j}=0$ otherwise. This is the trinomial distribution: the $n$ particles
are independent, and at time $t$ each is in one of three states, with probabilities
$c$ (neither inside nor outside --- that is, dead), $f$ (inside) and $g$ (outside).

\subsection*{Applied to the absorption models}

For absorption only, \cref{m:eq:examplePgf}, we have $c=0$, $f = 1-\ee^{-\alpha t}$
and $g = \ee^{-\alpha t}$. The factor $c^{\,n-i-j}$ then vanishes unless $i+j=n$,
and \cref{m:eq:stateProbSpec} gives
\begin{equation*}
p_{i,j}(t) = \binom{n}{i}\bigl(1-\ee^{-\alpha t}\bigr)^{i}\bigl(\ee^{-\alpha t}\bigr)^{j},
\qquad i+j=n,
\end{equation*}
so $X_t\sim\text{Binomial}\bigl(n,1-\ee^{-\alpha t}\bigr)$ and $Y_t = n-X_t$, as
anticipated in \cref{m:app:absonly}.

For absorption--death, \cref{m:eq:Gabsdeath}, the three functions are
\begin{equation*}
c(t) = 1-\frac{\mu}{\mu-\alpha}\ee^{-\alpha t} + \frac{\alpha}{\mu-\alpha}\ee^{-\mu t},
\qquad
f(t) = \frac{\alpha}{\mu-\alpha}\bigl(\ee^{-\alpha t}-\ee^{-\mu t}\bigr),
\qquad
g(t) = \ee^{-\alpha t},
\end{equation*}
which are precisely $p_{0,0}$, $p_{1,0}$ and $p_{0,1}$ of
\cref{m:eq:absdeathstates}.

% --- restore default section numbering for subsequent chapters --------------
\renewcommand{\thesection}{\thechapter.\arabic{section}}
